\section*{Thesis Summary} 
\todo[inline]{To finish.}
% Part I) Hypothesis-driven behavioral analysis for early autism screening

\textbf{\Cref{part:part_1}}: \textbf{Computational behavioral phenotyping for early autism screening.} This first part focuses on a closed-loop digital behavioral phenotyping approach to address the challenge of screening autism from video. The goal is to compute data-driven behavioral markers of autism using computer vision algorithms (CVA) for early autism screening. The first section reviews\\
We have been developing novel computational approaches for behavioral assessment of a wide range of autism symptoms with the goal of developing reliable, objective, quantitative, and scalable methods for behavioral measurement of symptoms of ASD and other developmental disorders \cite{CV1}.
Our goal is to design and empirically validated time-efficient methods, based on video analytics, for collecting, quantifying, and automatically interpreting behavioral data using relatively inexpensive, ubiquitous devices that clinicians, researchers, and caregivers can use to augment traditional methods of neurodevelopmental assessment. Caregiver autism screening questionnaires, while useful, have lower performance rates with caregivers who have lower education and are from minority ethnic/racial backgrounds; current methods often require relatively high levels of literacy and knowledge of child development \cite{Khowaja2014}. To what extent can digital tools based on direct behavioral observation assist in providing accurate and reliable assessments of autism behavioral symptoms? 

The tools we are developing for digital assessment methods are based on downloadable software applications for low-cost mobile devices (e.g., tablet or smartphone) that can automatically and objectively measure autism risk behaviors as indexed by the child’s visual attention \cite{CV_autism_1, team_article_2}, sensorimotor behaviors \cite{team_article_4}, and facial expressions \cite{Carpenter2020} in response to carefully-designed stimuli designed to elicit such behaviors. 


% Goals and questions:
% •	How can we capture early behavioral markers of autism from computer vision analysis? How can we combine them to estimate the diagnosis?
% •	How can we validate CVA estimates and quantify uncertainties?
% •	How to combine behavioral information? How to account for missing data? How can interpretability of the model’s predictions be recovered?


\begin{itemize}
    \item \textbf{\Cref{sec:chapter_1}}: \textbf{Definitions, study description, data collection, and features extraction}. This chapter introduces the study aim and data collection, presents the analyzed dataset and associated challenges.  TODO
    % •	Study description
    % •	Presentation of the data (app, clinical, and validity data)
    % •	Raw features extraction over time (head movement, gyro, head )
    % •	Evaluation and validation of the features extraction


    %------
    % the st the closed-loop digital behavioral phenotyping approach to address the challenge of screening autism from videos. The first section detail the study aims and data collection.  as a closed-loop behavioral phenotyping  and detail our approach  present the study description and data collection, the features extraction tackles the problem of searching for occurrences (repetitions) of predefined patterns within a single long time series. The first section reviews related work on time series similarity search, with a particular focus on exact methods that use lock-step distances. The algorithmic properties contributing to their efficiency are detailed. Building on these properties, the second section introduces a general framework for constructing distances invariant to user-defined sets of deformations while ensuring equivalent computation time complexity as state-of-the-art methods. The final section applies this framework to develop a distance invariant to amplitude scaling, offset shift, and linear trend. This distance proves valuable in cases where time series are affected by trend-induced deformations, as demonstrated experimentally.

    %presents the biomedical application motivating the development of the methods proposed in the subsequent chapters. 
    %Shortly, an enzyme plays an important role in regulating muscle activity and signal transmission within the nervous system. Some drugs inhibit the enzyme's action, which has severe consequences, notably on respiration, which are not yet fully understood. To investigate inhibition consequences, the respiration of mice with different genotypes exposed to inhibitors is monitored by plethysmography. 
    % The first section presents the monitoring equipment material (plethysmogram), and highlights the limitations of existing methods for analyzing such signals. Additionally, an algorithm for segmenting plethysmography signals into datasets of respiratory cycles (inspiration and expiration) is presented. 
    % The second section outlines the biological context and experimental protocol.

    \item \textbf{\Cref{sec:chapter_2}: Behavioral markers estimation from multimodal data, validation, and uncertainty quantification}. 
        \begin{itemize}
            \item In \Cref{sec:chapter_2_nc}, we propose two algorithms based on computer vision and audio signal processing to detect the child's head turning and name calls utterances accurately, and report findings related to diagnosis-groups behavioral differences captured by this experiment. In \Cref{sec:chapter_1_nc_motivation} we present the related work and evidence supporting the addition of this component. 

            % [FROM NM] While children were watching the movies, their name was called three times by an examiner standing behind them. The analysis of the video audio allowed us to automatically detect the instant in which the participant's name was called, while the dynamics of the head orientation (pitch, yaw, and roll) were used to automatically detect if and when the child responded (oriented) toward the person calling their name. Computed vision and audio algorithms were originally proposed in \cite{CV_autism_1} and improved in \Cref{journal_2} (see \Cref{sec:chapter_2}).  

            
            
            In this chapter we present a computer vision–based digital phenotyping tool to quantify toddlers' behavioral responses to name-calling during pediatric visits. Using app-delivered video stimuli and automated head-turn detection on smartphones/tablets, we demonstrated that children with ASD showed significantly reduced response frequency and increased latency, replicating human coding and enabling early symptom detection through fine-grained spatiotemporal behavioral analysis.

            We develop and apply fine-grained mathematical models to capture key aspects of gameplay, such as a refined measure of finger's accuracy across time using a 3D-spatiotemporal modeling of the game (see \Cref{sec:chapter_2_accuracy_computation}), and a novel proxy for the force applied to the screen estimated from the device's supporting tripod dynamics during touches. These models enable the extraction of interpretable motor phenotypes at sub-second resolution, moving beyond coarse behavioral metrics and offering insight into subtle developmental signatures.

            % used in intro of Chap 2-2:  Finally, we propose a mathematical modeling of the child's finger-touchscreen interactions to represent fine-grained aspects of the gameplay sessions. This enabled to extract and evaluate a series of novel visual-motor features from raw gameplay data, and explore associations with children’s cognitive, language, and motor abilities, as well as autism-related behaviors.




            \item Increasing evidence suggests that early motor impairments are a common feature of autism. This chapter presents an engaging and scalable assessment of the visual-motor abilities based on a bubble-popping game administered on a tablet. Participants were 233 children ranging from 1.5 to 10 years of age (147 neurotypical children and 86 children diagnosed with ASD, of which 32 were also diagnosed with co-occurring attention-deficit/hyperactivity disorder). In this chapter, we introduce a quantitative modeling of game sessions to extract several game-based touch features, which were compared across autistic, autistic+ADHD, and neurotypical participants. Among younger children (1.5-3 years), autistic children popped the bubbles at a lower rate, and their ability to touch the bubble’s center was less accurate compared to neurotypical children. When they popped a bubble, their finger lingered for a longer period, and they showed more variability in their performance. In older children (3-10-years), consistent with previous research, the presence of co-occurring ADHD was associated with greater motor impairment, reflected in lower accuracy and more variable performance. Several motor features were correlated with standardized assessments of fine motor and cognitive abilities, as evaluated by an independent clinical assessment. 
        \end{itemize}

    \item \textbf{\Cref{sec:chapter_3}}: \textbf{Multimodal machine learning for interpretable digital behavioral phenotyping and application to early autism screening}. 
    
    The last chapter describes all of the 23 autism-related digital phenotypes assessed via the app, and presents our multimodal data processing system $\mathcal{S}$ used to differentiate autistic toddlers from their neurotypical counterparts. We evaluated its accuracy for identification of autism in a large sample of 475 toddler-age children.  The resulting algorithm achieved high diagnostic accuracy (AUC = 0.90, sensitivity = 87.8\%, specificity = 80.8\%) and robust performance across sex, race, and ethnicity. We conclude that quantitative, objective, and scalable digital phenotyping offers promise in increasing the accuracy of autism screening and reducing disparities in access to diagnosis and intervention.

    \item We introduce a prediction confidence index, quantifying the model prediction certainty. This score is derived from the consistency of predictions across 1,000 stratified resampling experiments, where each participant’s classification is recorded. We show that this score aligns with the consistency of SHAP-based feature attributions, with most conclusive predictions (respectively least conclusive) are associated with uniformly distributed SHAP attributions towards autistic or neurotypical-related behaviors (respectively more heterogeneous) .
    % •	CVA app features presentation 
    % •	Analysis of missing data and imputation techniques 
    % •	Benchmark across learning algorithms and XGBoost presentation 
    % •	Tabular data classification for early autism screening, class-imbalance handling, appropriate performance evaluation, uncertainties quantification
    % •	Behavioral phenotyping using missing data-aware SHAP values to enhance predictions’ interpretability

\end{itemize} 


\textbf{\Cref{part:part_2}}: \textbf{Data-driven estimation of egocentric vision ethograms with application to the ecological assessment of cognitive and behavioral executive dysfunction.} This part focuses on the development of data processing methodologies to characterize behavioral alterations induced by sensorimotor disorders and dysexecutive syndromes. The first chapter presents the traditional and computational assessment methods, the study design, data collection, and features extraction from the raw multimodal data. In the second chapter, we present a procedure to estimate a set of reliable, interpretable and unequivocal egocentric vision atoms to describe the participant's egocentric vision trajectory. The egocentric vision ethograms of healthy participants are then used to estimate a barycentric symbolic representation for the prescribed task, enabling to investigate within- and inter-group differences in behavioral phenotypes in \Cref{sec:chapter_6} 

% Goals and questions:
% •	Quantify the participant’s behavior when performing a prescribed task. 
% •	Can we temporally segment a video (what are segments, actions, activities, etc)? 
% •	Can we estimate a common, interpretable, and exhaustive ego-vision vocabulary to describe the task from an egocentric viewpoint?  
% •	How can we use the reconstructed ethogram or symbolic representation to compare behaviors across diagnosis groups?


\begin{itemize}
    \item \textbf{\Cref{sec:chapter_4}}: \textbf{Definitions, study description, data collection, and features extraction}.     
    %\textbf{\Cref{sec:chapter_1}: Definitions, study description, data collection, and first-order features extraction}
   TODO
        •	Description of the Smartflat study and grounding in neuro-ergonomy and neurology
        •	Multimodal and multi-view, longitudinal data collection (data presentation)
        •	Extraction of multiple modalities
        o	Spatiotemporal video sequence
        o	Speech recognition and diarization
        o	Hands and body landmarks
        o	Gaze and head kinematic data 
        

    %\todo[inline]{I think the ethogram introduced as a tool in ethology (at fist to study and quantify animal's behavior) is the appropriate name for bridging more with the clinical interface side. At the same time, perhaps we should stick more to mathematical terminology and use terms like "symbolic representation" instead of ethogram, and be consistent throughout the manuscript to only refer to ethogram and associated tools when using them as is. Any input or decisive feedback is welcome, you all have more experience :-) -> Self-solve: explain this in detial when introducing the ethogram in \Cref{sec:smartflat_introduction} and \Cref{sec:symbolization}.}

    %Specifically, this global scale task is addressed by considering groups of elastic deformations that can either be restricted or quantified and embedded. 

    
    \item \textbf{\Cref{sec:chapter_5}}: \textbf{Recursive ego-vision syllabes prototyping using latent subspace exploration with human feedback}. TODO

    % The mapping $\Phi$ aggregates the K first-stage centroids into $G$ higher-level prototypes, in a manner consistent with the hierarchical organization typically found in ethograms.

    %  This chapter presents a novel unsupervised baseline method for analyzing egocentric vision data. The method uses a recursive exploration of the latent space with human feedback to prototype ego-vision syllables, which are then used to estimate an ethogram of the participant's behavior during a cooking task. The results show that the proposed method is effective in estimating the ethogram and can be used to compare behaviors across diagnosis groups. 
    

    % %     \item \textbf{\Cref{sec:change_point_detection}} \textbf{Temporal segmentation and activity recognition}.     

    % Ego-vision spatiotemporal token prototyping: an effective and constructive exploration of a cooking-task latent subspace with application to ethogram estimation
    
    % •	What is a spatiotemporal token prototype? What does it mean and how can we qualify and quantify it? Motivations for using a prototypes-based approach
    % •	Recursive exploration of the latent space with human feedback
    % •	Hierarchical clustering of prototypes for higher-semantic lower-dimensional visual tokens.
    % •	Evaluation and ablation study 

    % •	Multivariate time-series approach with change-points detection. 
    % •	Clustering within or across participants for activity recognition
    % •	Performance metrics and state-of-the-art
    
    %presents, in the first section, a novel unsupervised baseline method for analyzing plethysmography signals. This method uses a DTW-based clustering algorithm to learn a symbolic embedding of respiratory cycles. The symbolization of plethysmography signals results in sequences of shape-based symbols. The following result and discussion illustrate, in particular, the interpretability of the method by presenting correspondences between symbols and physiological functions. Among several discoveries, the symbolic representations have highlighted genotype-dependent respiration modalities and a heterogeneous physiological response to inhibitor exposures. 

    \item \textbf{\Cref{sec:chapter_6}}:  \textbf{Symbolic representation and barycenter estimation with applications to dysexecutive syndromes behavioral and cognitive assessment.}.     
    
    % •	Application of the ego-vision vocabulary from symbolic representation, with quantitative and qualitative validation 
    % •	Custom barycenter estimation using Shape-based Dynamic Time Warping Barycenter Averaging (ShapeDBA) 
    % •	Application of the symbolization and barycenters for the comparison across severity-stratified diagnosis groups, and interpretable behavioral phenotyping 
    % •	Opportunities with additional multimodality (with gaze and oculomotor behavior) and multi-view (three additional GoPro views). 
    

\end{itemize}



\textbf{\Cref{part:part_3}}: \textbf{Cardiovascular health monitoring from wearable physiological data.} focuses on unsupervised 

\begin{itemize}
    \item \textbf{\Cref{sec:chapter_7}}: \textbf{Time-varying Representations of Longitudinal Biosignals using Self-supervised Learning.}.  
    This chapter investigates how integrating timestamp information into contrastive self-supervised learning improves the representation of wearable biosignals for health monitoring. By designing time-aware positive pair sampling strategies — leveraging proximity in date and time of day — we enforce physiological consistency in learned embeddings. We show that such representations better capture circadian patterns and improve the longitudinal prediction of key health metrics from Photoplethysmography (PPG) alone. Our results emphasize the value of temporally sensitive training and evaluation protocols for biosignal-based health assessment.   
    
        % •    Introduction to the Apple Heart Study and the associated dataset
        % •    Self-supervised learning for time-varying representations of longitudinal biosignals
        % •    Evaluation of the learned representations on downstream tasks
        % •    Comparison with state-of-the-art methods

\end{itemize}